\documentclass[11pt]{article}
\usepackage[margin=1in]{geometry}
\usepackage[T1]{fontenc}\usepackage[utf8]{inputenc}\usepackage{lmodern}
\usepackage{amsmath,amssymb,amsthm,mathtools,booktabs,microtype}
\usepackage[colorlinks=true,linkcolor=blue,urlcolor=blue,citecolor=blue]{hyperref}
\hypersetup{pdftitle={Depth bounds for finite Gaussian Bayesian learning},pdfauthor={ChatGPT-6 Astra; Claude Opus 5.5; Kaarel Hänni; Dmitry Vaintrob}}
\newtheorem{theorem}{Theorem}[section]\newtheorem{lemma}[theorem]{Lemma}\newtheorem{proposition}[theorem]{Proposition}\newtheorem{corollary}[theorem]{Corollary}
\newcommand{\E}{\mathbb E}\newcommand{\R}{\mathbb R}\newcommand{\KL}{\operatorname{KL}}\newcommand{\op}{\mathrm{op}}
\title{Depth bounds for finite\\Gaussian Bayesian learning}
\author{ChatGPT-6 Astra \and Claude Opus 5.5 \and Kaarel H\"anni \and Dmitry Vaintrob\thanks{Authors are listed in alphabetical order.}}
\date{17 November 2026}
\begin{document}\maketitle
\begin{abstract}
We prove exponential-depth sample bounds for finite Gaussian Bayesian neural networks without a robustness restriction or a mixture over architectures. At fixed input dimension, comparison width, coefficient cap, accuracy, and confidence, $\exp(cH)$ samples suffice to compete with a depth-$H$ comparison network in squared loss. For bounded centered hidden activations such as tanh, this holds at every prior scale $0<\tau\le1$, including canonical mean field, with raw Gaussian likelihood and the ordinary raw posterior mean. The new comparison construction uses geometrically increasing groups of active clones: only first-layer noise is suppressed, and its normalized relative-entropy cost is linear in depth. A second all-scale theorem covers ReLU hidden layers with a bounded final readout, such as clipping or tanh, used in both likelihood and prediction. At every fixed power-law exponent $\gamma>1$, the raw-output theorem also permits bounded noncentered activations and has a growing-dimension extension. All theorems hold simultaneously at every learner width above the comparison width, for arbitrary bounded IID labels. For bounded noncentered activations at mean field, we give weaker depth bounds; a smooth-activation bound is polynomial in dimension at fixed depth. The constants may depend on the fixed comparison resources. These are exponential-depth upper bounds, not polynomial-depth guarantees.
\end{abstract}

\section{Statement and setup}
A sample bound that is polynomial in comparison width at each fixed depth can still grow much faster than exponentially with depth. This paper gives exponential-depth bounds for several Gaussian Bayesian models, including raw-output tanh at canonical mean field. We keep the learner's prior fixed and allow its likelihood calibration to depend on width and the requested accuracy. The models and scope are distinguished below:
\begin{itemize}
\item Bounded centered hidden activations and raw outputs: every scale $0<\tau\le1$, at fixed input dimension (Theorem~\ref{thm:geometric-depth}).
\item ReLU hidden layers with a bounded final readout in both likelihood and prediction: every scale $0<\tau\le1$, at fixed dimension (Theorem~\ref{thm:relu-all-scale-depth}).
\item Bounded hidden activations, allowing a nonzero value at zero, and raw outputs: every fixed $\gamma>1$, with an extension to dimensions up to depth under a row-norm comparison bound (Theorem~\ref{thm:main} and Corollary~\ref{cor:joint-dimension}).
\end{itemize}

Fix input dimension $d\ge1$, comparison width $C\ge1$, coefficient cap $B\ge1$, accuracy $0<\epsilon<1$, confidence $0<\delta<1$, and, for the first theorem, exponent $\gamma>1$. Let $\phi:\R\to[-1,1]$ be globally one-Lipschitz. No smoothness or condition $\phi(0)=0$ is assumed. At hidden depth $H\ge2$ and width $N$, the architecture is
\[
 h^1(x)=\phi(W^1\bar x),\quad
 h^\ell(x)=\phi(W^\ell h^{\ell-1}(x))\ (2\le\ell\le H),\quad
 z_w(x)=v^\top h^H(x),\qquad \bar x=(1,x).
\]
There are no other bias parameters. Let $\mathcal G_{C,B,H}$ be the same-depth width-$C$ comparison class with every coefficient bounded by $B$. Its outputs have absolute value at most $U=CB$.

Inputs and labels have any joint law $D$ on $[-1,1]^d\times[-1,1]$, and the sample is IID. Write $R(f)=\E_D(f(X)-Y)^2$ and $\widehat R_S(f)=n^{-1}\sum_i(f(X_i)-Y_i)^2$.
For any $0<\tau\le1$ and $\omega=N/\tau^2$, independent centered Gaussian prior variances are
\[
 \tau^2/d\text{ in }W^1,\qquad
 \tau^2/N\text{ in internal matrices},\qquad
 \tau^2/N^2\text{ in }v.
\]
A Gaussian likelihood of variance $n\lambda/\omega$ gives the actual finite posterior
\begin{equation}
 d\Pi_{N,H,S}(w)=Z^{-1}e^{-\omega\widehat R_S(z_w)/(2\lambda)}d\mu_{N,H}(w).
 \label{eq:posterior}
\end{equation}

\begin{theorem}[Exponential depth suffices above mean field]
\label{thm:main}
For $\tau^2=N^{1-\gamma}$ with the fixed $\gamma>1$, there are constants $a,c>0$ and $0<\ell_+,\ell_-<1$, depending only on $d,C,B,\epsilon,\delta,\gamma$, such that the following holds for every $H\ge2$. Define
\[
 N_*(H)=\left\lceil e^{2aH/(\gamma-1)}\right\rceil,
 \qquad
 \lambda_{H,N}=
 \begin{cases}
 \ell_-/H^2,&N<N_*(H),\\
 \ell_+/H,&N\ge N_*(H).
 \end{cases}
\]
If $n\ge e^{cH}$, then with probability at least $1-\delta$, simultaneously for every $N\ge C$,
\begin{equation}
 \int R(z_w)d\Pi_{N,H,S}
 \le\inf_{g\in\mathcal G_{C,B,H}}R(g)+\epsilon.
 \label{eq:main}
\end{equation}
The same bound holds for the ordinary raw posterior mean. The constants can be chosen uniformly over all activations satisfying the stated range and Lipschitz bounds.
\end{theorem}

\begin{corollary}[Dimension may grow with depth]
\label{cor:joint-dimension}
Suppose the comparison network's first-layer rows, including their bias coordinate, have Euclidean norm at most $B$, and later coefficients have absolute value at most $B$. For every $1\le d\le H$, Theorem~\ref{thm:main} holds with the constants $a,c,\ell_\pm$ independent of $d$. They depend only on $C,B,\epsilon,\delta,\gamma$. In particular the same $e^{cH}$ sample bound applies to arbitrary input laws on the unit sphere in dimension at most $H$.
\end{corollary}

The theorem allows arbitrary bounded label noise and approximation error. It does not change the Gaussian prior as the sample count grows. The prior exponent is fixed strictly above one; constants may diverge as $\gamma\downarrow1$. At mean field, Theorem~\ref{thm:geometric-depth} proves the same fixed-dimension exponential-depth conclusion when $\phi(0)=0$. Section~\ref{sec:all-scales} gives weaker bounds without that centering assumption. None of these statements supplies a bound polynomial jointly in $d,C,B,H$.

\section{The comparison cost with its width factor retained}
The quadratic form in the prior density is
\[
 Q_N(w)=\frac dN\|W^1\|_F^2+
          \sum_{\ell=2}^H\|W^\ell\|_F^2+N\|v\|^2,
 \qquad d\mu\propto e^{-\omega Q_N/2}\,dw.
\]
Partition learner neurons into $C$ nearly equal groups and clone the comparison network, dividing outgoing edges by source-group sizes. The resulting network computes $g$ exactly and has energy
\[
 E_H=d(d+1)B^2+2HC^2B^2\le c_0H,
\]
where constants $c_i$ in this section depend only on $d,C,B$.

\begin{lemma}[A width-dependent Gaussian neighborhood cost]
\label{lem:teacher}
Let $0<\lambda\le1$, $\rho=\lambda/[4(U+1)]$, and
\[
 D_H=d+2+(H-1)C,\quad
 P_H=1+2CB(d+H+2)(2+3CB+8)^{H-1},\quad
 \eta=\min\{1/2,\rho/(2P_H)\}.
\]
The prior mass of $\|z_w-g\|_\infty<\rho$ is at least $e^{-\omega A_{H,N}(\lambda)}$, where
\begin{equation}
 A_{H,N}(\lambda)=\frac{E_H}{2}
     +\tau^2D_H\log\frac8{3\eta}
     +\frac{(H-1)\log2}{\omega}.
 \label{eq:exact-comparison-cost}
\end{equation}
Consequently the posterior empirical objective is at most
\begin{equation}
 \int\widehat R_S(z_w)d\Pi+\frac{2\lambda}{\omega}\KL(\Pi\Vert\mu)
 \le\widehat R_S(g)+\lambda+2\lambda A_{H,N}(\lambda).
 \label{eq:comparison}
\end{equation}
\end{lemma}
\begin{proof}
Use standard Gaussian coordinates. Require absolute value at most $\eta$ for all first-layer and readout coordinates and for each internal row's $C$ coordinates in the group-constant subspace. There are $ND_H$ independent constraints, with probability at least $(3\eta/8)^{ND_H}$.

A centered Gaussian truncated to a symmetric interval remains unit-subgaussian. Orthogonal unconstrained coordinates are still independent standard Gaussians. Conditional on the cylinder, a pair of $1/4$ sphere nets therefore gives probability at least $1/2$ that each internal standardized matrix has operator norm at most $8\sqrt N$. Different internal blocks remain independent. Their joint conditional probability is at least $2^{-(H-1)}$; no depth-dependent operator cutoff is needed.

Translate this symmetric event by the standardized clone. Normalized hidden-state errors satisfy
$e_1\le\eta(d+1)$ and
$e_\ell\le(\sqrt2CB+8)e_{\ell-1}+\eta\sqrt C$.
The readout perturbation then gives uniform output error at most $\eta P_H\le\rho/2$. For a standard Gaussian measure and a symmetric event $A$, translation by $b$ decreases its mass by at most $e^{-\|b\|^2/2}$: the remaining density ratio has conditional expectation at least one by Jensen. Here $\|b\|^2=\omega Q_N(w_*)\le\omega E_H$. This proves \eqref{eq:exact-comparison-cost}. Conditioning the prior on the output neighborhood gives entropy at most $\omega A_{H,N}$ and empirical loss at most $\widehat R(g)+2\rho(U+1)+\rho^2\le\widehat R(g)+\lambda$. The posterior variational identity proves \eqref{eq:comparison}.
\end{proof}

For either $\lambda=\ell/H$ or $\lambda=\ell/H^2$,
\begin{equation}
 A_{H,N}(\lambda)\le c_1H+
      c_2\tau^2H^2[1+\log(1/\ell)].
 \label{eq:teacher-depth}
\end{equation}
This follows from $D_H=O(H)$, $\log P_H=O(H)$, and $\log H\le H$. At every width it is $O(H^2)$; when $\tau\le e^{-aH}$ with $a\ge1$, it is $O(H)$. The constants in these statements depend on $\log(1/\ell)$. Choose $\ell_\pm$ sufficiently small using $\ell[1+\log(1/\ell)]\to0$, so the comparison penalties in \eqref{eq:comparison} are at most a prescribed small fraction of $\epsilon$ in their respective width regimes.

\section{Spectral energy and layer products}
Retain from each internal matrix only singular values greater than $8\tau$, obtaining $T^\ell$. Leave the first layer and readout unchanged. Let $\widetilde z$ be this truncated network and write
\[
 Q=\frac dN\|W^1\|_F^2+
       \sum_{\ell=2}^H\|T^\ell\|_F^2+N\|v\|^2.
\]

\begin{lemma}[Gaussian moment and spectral approximation]
\label{lem:spectral}
There is a constant $c_d$ depending only on dimension such that
\begin{equation}
 \log\E_\mu e^{\omega Q/16}\le c_dN+(H-1)\log(3/2).
 \label{eq:energy-mgf}
\end{equation}
Also $|z|,|\widetilde z|\le\sqrt Q$ and
\begin{equation}
 \|z-\widetilde z\|_\infty
 \le8\tau(H-1)e^{Q/(2e)}.
 \label{eq:spectral-depth}
\end{equation}
\end{lemma}
\begin{proof}
The first-layer and readout terms, multiplied by $\omega$, are chi-square variables with $N(d+1)$ and $N$ degrees of freedom. For an internal matrix, the operator-norm ball of radius $4\tau$ has prior probability at least one half. Its squared Frobenius distance is $\sum_j(s_j-4\tau)_+^2$, and the retained energy is at most four times this distance squared. Gaussian isoperimetry gives retained-energy tail at most $\tfrac12e^{-\omega t/8}$. Integration bounds its exponential moment at $\omega/16$ by $3/2$. Independence proves \eqref{eq:energy-mgf}.

For the approximation, telescope replacements so that the prefix is original and the suffix is truncated. Bounded activations make the incoming hidden norm at every replacement at most $\sqrt N$. The layer-$\ell$ contribution is at most
$8\tau\sqrt N\|v\|\prod_{j>\ell}\|T^j\|_\op$.
A product of $k$ nonnegative factors whose squared sum is at most $Q$ is at most $(Q/k)^{k/2}$ by AM--GM. Moreover
\[
 (Q/k)^{k/2}\le e^{Q/(2e)}
\]
for every positive $k$, by maximizing $k\log(Q/k)$. Summing the $H-1$ replacements proves \eqref{eq:spectral-depth}. The output envelope follows directly from bounded hidden activations and Cauchy--Schwarz.
\end{proof}

\section{Capacity with exponential constants}
We state explicitly the Gaussian-process ingredient that avoids a factor such as $(\log n)^H$. For a symmetric compact $F\subset[-1,1]^n$ containing zero, let $G(F)=n^{-1}\E_g\sup_{f\in F}\sum_i g_i f_i$, and let $\mathcal L$ be scalar one-Lipschitz functions on $[-1,1]$ vanishing at zero. Then
\begin{equation}
 G(\mathcal L\circ F)\le cG(F)+c/\sqrt n
 \label{eq:scalar-link}
\end{equation}
for an absolute $c$.

Here is a proof from standard generic chaining \cite{talagrand}. The scalar class has uniform entropy $\log\mathcal N(a,\mathcal L,\|\cdot\|_\infty)\le c/a$, so its $\gamma_2$ functional is bounded by an absolute constant. Equip $F$ with $d_n(f,f')=\|f-f'\|_2/\sqrt n$. Product admissible nets, shifted by one level, give
$\gamma_2(F\times\mathcal L,d_n+\|\cdot\|_\infty)\le c[\gamma_2(F,d_n)+1]$.
Composition is one-Lipschitz for this sum metric. The generic-chaining upper bound and the majorizing-measure lower bound for the Gaussian process on $F$, whose canonical metric is $d_n/\sqrt n$, prove \eqref{eq:scalar-link}.

Start with the absolute convex hull $\mathcal A_0$ of the constant function and input coordinates. Recursively take the absolute convex hull of the constant and $\mathcal L\circ\mathcal A_{\ell-1}$. These classes have envelope one. Absolute convexification preserves Gaussian width; adding a constant costs at most $1/\sqrt n$. Thus their Rademacher complexities satisfy
\begin{equation}
 \mathfrak R_n(\mathcal A_H)
 \le\frac{c^{H+1}[1+\sqrt{\log(2d+2)}]}{\sqrt n}.
 \label{eq:atoms}
\end{equation}
We used $\mathfrak R(F)\le\sqrt{\pi/2}\,G(F)$, obtained by conditioning Gaussian signs and applying Jensen to their magnitudes.

For a given network, let $\beta=|\phi(0)|\le1$ and assign neuron budgets
$r_j^1=\beta+\sum_i|W^1_{ji}|$ and
$r_j^\ell=\beta+\sum_i|T^\ell_{ji}|r_i^{\ell-1}$.
At preactivation budget $t>0$, the normalized link
$[\phi(tu)-\phi(0)]/t$ lies in $\mathcal L$.
Induction therefore places the network output in $P\mathcal A_H$, where $P=\sum_j|v_j|r_j^H$. Cauchy--Schwarz and expansion of the budget recurrence give
\begin{equation}
 P\le\sqrt2\sum_{k=1}^{H+1}(Q/k)^{k/2}
 \le\sqrt2(H+1)e^{Q/(2e)}.
 \label{eq:path}
\end{equation}
The factor $\sqrt2$ accounts for the first-layer bias coordinate, since $(d+1)/d\le2$. Each expanded term contains a subset of layer norms and the readout norm, so the same total-energy AM--GM bound applies.

Consequently, at $Q\le TH$ for any fixed $T$, the union of output classes over all widths has envelope $\sqrt{TH}$ and worst-case empirical Rademacher complexity at most $e^{c_{T,d}H}/\sqrt n$. A standard smooth-square-loss bound \cite{sst} gives, on one event, for every function in this class with empirical risk at most sixteen,
\[
 \begin{aligned}
 R(f)-\widehat R(f)\le c\bigg[&
 \sqrt{\widehat R(f)}\left(\log^{3/2}(2n)\mathfrak R_n+
       \sqrt{(TH+1)\log(c/\delta)/n}\right)\\
 &+\log^3(2n)\mathfrak R_n^2+(TH+1)\log(c/\delta)/n\bigg].
 \end{aligned}
\]
For fixed $T,d,\epsilon,\delta$, this is at most $\epsilon/32$ whenever $n\ge e^{c'H}$ for sufficiently large $c'$. The remaining polynomial powers of $H$ and $\log n$ are absorbed by enlarging $c'$. Rational finite-parameter approximations at each width supply separability for the union.

\section{Posterior integration at large widths}
Choose $\ell_+$ as after \eqref{eq:teacher-depth}, and put $\lambda=\ell_+/H$. Since the centered readout gives $\E_\mu\widehat R\le2$, Jensen yields $Z\ge e^{-\omega/\lambda}$. Also $r=N\|v\|^2$ has $\E r\le1$, $\E r^2\le3$, and $R(z)\le2r+2$, so $\E_\mu R(z)^2\le24$.

First choose a sufficiently large fixed $T$. Equation~\eqref{eq:energy-mgf} and Cauchy--Schwarz bound the posterior energy-tail risk by
\[
 \int R(z)\mathbf1_{Q>TH}d\Pi
 \le\sqrt{24}\exp\left[\frac\omega\lambda+
       \frac{c_dN+(H-1)\log(3/2)}2-\frac{\omega TH}{32}\right].
\]
Since $N\le\omega$ and $H\ge2$, this is at most a prescribed small fraction of $\epsilon$ when $T$ is sufficiently large depending on the fixed resources. The high-empirical-risk tail is at most
$2e^{-3\omega/\lambda}$ on $\widehat R(z)>8$.

Next choose $a\ge1$ so that $8H e^{-aH+TH/(2e)}$ is sufficiently small compared with $\epsilon$ for every $H\ge2$. At $N\ge N_*(H)$, $\tau\le e^{-aH}$, so \eqref{eq:spectral-depth} gives that uniform approximation accuracy on $Q\le TH$. The preceding statistical event bounds the truncated network's risk. Minkowski's inequality for empirical and population square-loss norms transfers the estimate back to $z$, yielding
$R(z)\le\widehat R(z)+\epsilon/4$ on $Q\le TH$, $\widehat R(z)\le8$, after choosing the intermediate constants more conservatively if needed.

Integrate this inequality, add both tail contributions, and use \eqref{eq:comparison}. By the choice of $\ell_+$ the result is
$\int R(z)d\Pi\le\widehat R(g)+\epsilon/2$.
The event is common to all large widths because the energy-class statistical bound includes their union.

\section{The remaining widths and the final comparison}
We give the finite-width bound used here to keep the argument independent of other manuscripts. For $0<\lambda\le1$, proof-clip outputs to $[-M,M]$ with $M^2=128/\lambda$. The square loss is bounded by $b\le512/\lambda$. Put $\beta=\omega/(2\lambda)$ and $s=\beta b/n$. For $s\le1$, the elementary exponential PAC--Bayes inequality gives, on one event, every law $Q$ the bound
\[
 \int R(z^{[M]})dQ\le\frac{s}{1-e^{-s}}
 \left[\int\widehat R(z^{[M]})dQ+
          \frac{\KL(Q\Vert\mu)+\log(1/\delta_N)}\beta\right].
\]
It follows by convexity of $e^{-s\ell/b}$ for $\ell\in[0,b]$, independence, integration over the prior, Markov's inequality, and entropy duality. The actual posterior's clipped objective is at most its raw objective, which is at most two by comparison with its own prior. Since $s/(1-e^{-s})\le1+s$, the comparison lemma gives
\begin{align}
 \int R(z)d\Pi\le\;&\widehat R(g)+\lambda+2\lambda A_{H,N}(\lambda)
       +\frac{4\lambda\log(1/\delta_N)}\omega
       +\frac{512\omega}{n\lambda^2}
       +\sqrt{24}e^{-7\omega/\lambda}.
 \label{eq:finite-width}
\end{align}
The last term restores raw outputs: the Gaussian readout tail gives
$\mu(r>M^2)\le e^{-\omega M^2/8}$, and Cauchy--Schwarz with the evidence lower bound bounds its posterior risk by that term. The calibration condition $256\omega/(n\lambda^2)\le1$ ensures $s\le1$.

For $N<N_*(H)$ take $\lambda=\ell_-/H^2$ and $\delta_N$ proportional to $\delta/N^2$. The confidence ratio is uniformly bounded:
$\log(1/\delta_N)/\omega\le c[\log(1/\delta)+1]$.
The teacher cost is $O(H^2)$, with its $\log(1/\ell_-)$ dependence already accounted for. Choose $\ell_-$ small enough that its comparison, confidence, and raw-tail terms are small fractions of $\epsilon$. Finally
\[
 \omega=N^\gamma\le\exp(c_{a,\gamma}H),\qquad
 \lambda^{-2}=H^4/\ell_-^2.
\]
Thus $n\ge e^{cH}$ makes the sample term and its calibration condition valid uniformly at these widths. The sum of the width failure probabilities is finite.

Allocate the remaining confidence to one fixed teacher's scalar Bernstein event. Since $|g|\le CB$, a sample count independent of depth already makes $\widehat R(g)\le R(g)+\epsilon/2$; it is absorbed by $e^{cH}$. The comparison class is compact, so a population-risk minimizer exists, and it has risk at most one because zero readout is allowed. Apply the proof to this fixed minimizer. Intersecting the large-width, finite-width, and teacher events proves Theorem~\ref{thm:main}. Each posterior has finite readout moments, and Jensen proves its ordinary-mean assertion.

The order of choices is important: choose $\ell_\pm$ using $\ell(1+\log(1/\ell))\to0$, then $T$, then $a$, then the sample exponent $c$. No constant selected later is needed to justify an earlier choice.

\begin{proof}[Proof of Corollary~\ref{cor:joint-dimension}]
The row-norm bound improves the first clone energy to $dB^2$, so $E_H\le dB^2+2HC^2B^2=O(H)$ uniformly for $d\le H$. Also $D_H=O(H)$ and $\log P_H=O(H)$ uniformly there. Hence the same comparison-cost choices are dimension-independent. In the energy-tail proof $c_d=O(d)$, so $c_dN/(\omega H)=O(1)$ when $d\le H$. The only dimension factor in the atom complexity is $1+\sqrt{\log(2d+2)}$, which is absorbed into $e^{cH}$ uniformly. The remaining finite-width and teacher concentration arguments have no dimension-dependent constants. The cube statement includes the unit sphere as a subset.
\end{proof}

\section{Limits of the result}
The bound is exponential in depth at fixed comparison resources and fixed accuracy. Its likelihood calibration uses those declared resource bounds, but not the unknown data law or comparison minimizer. It does not imply efficient posterior computation.

At mean field, $\tau=1$ regardless of width. The large-width spectral approximation used here therefore never becomes accurate. The next section gives a weaker all-scale upper bound by a different argument. Whether its depth dependence can be improved to exponential remains open. Unbounded activations also require different control of the raw-output finite-width tail; final saturation of both likelihood and prediction is a separate model.

\section{Canonical mean field with bounded centered activations}
\label{sec:geometric-clones}
For bounded activations satisfying $\phi(0)=0$, a different comparison network removes the strictly-above-mean-field assumption from the fixed-dimension exponential-depth theorem. Its active width grows geometrically through the hidden layers. This is a comparison distribution inside the same dense Gaussian prior; the learner's architecture and prior support remain unchanged.

\begin{theorem}[All scales, raw outputs, bounded centered activation]
\label{thm:geometric-depth}
Fix $d,C,B,\epsilon,\delta$. Let $\phi:\mathbb R\to[-1,1]$ be one-Lipschitz with $\phi(0)=0$. There are constants $a,c>0$ and $0<\ell_+,\ell_-<1$, depending only on the fixed resources, such that for every $H\ge2$, the following calibration has the stated guarantee:
\[
 N_*(H)=\lceil e^{aH}\rceil,\qquad
 \lambda_{H,N}=
 \begin{cases}\ell_-/H^2,&N<N_*(H),\\
 \ell_+/H,&N\ge N_*(H).
 \end{cases}
\]
If $n\ge e^{cH}$, one sample event of probability at least $1-\delta$ satisfies \eqref{eq:main} simultaneously for every $N\ge C$ and every $0<\tau\le1$. The posterior uses the raw Gaussian likelihood \eqref{eq:posterior}; the conclusion bounds posterior-average risk and the risk of the ordinary raw posterior mean. In particular canonical mean field $\tau=1$, with internal variances $1/N$, is included.
\end{theorem}

This theorem covers tanh without a robustness restriction. The dependence on dimension and comparison resources is held fixed in the exponential-depth statement. The centering assumption is used to keep unused neurons close to zero; it is not needed for the earlier strictly-above-mean-field theorem.

\subsection{Geometrically increasing clone groups}
Write the width-$C$ comparison matrices and readout as $V^1,\ldots,V^H,b$, and put $A=CB$. Choose an integer
$\alpha\ge2A^2+4$. A constant $0<p\le1/4$ will be fixed by the requested approximation accuracy, independently of $H,N,\tau$. Whenever
\[
 N\ge\frac{2C\alpha^{H-1}}p,
\]
define
\[
 k_1=\left\lfloor\frac{pN}{C\alpha^{H-1}}\right\rfloor,
 \qquad k_l=k_1\alpha^{l-1},\qquad
 m_l=Ck_l,\qquad p_l=m_l/N.
\]
Then $p/2\le p_H\le p$, $p_l=p_H\alpha^{l-H}$, and every layer has at most $N$ active neurons. Each comparison neuron has $k_l$ active clones in layer $l$. Give first-layer clones the comparison incoming row, connect every layer-$l$ active clone to every clone of its comparison parents with weights $V^l_{ij}/k_{l-1}$, and give each final clone readout coefficient $b_j/k_H$. All other mean weights are zero. The mean network computes the comparison function exactly: inactive mean neurons have output $\phi(0)=0$.

Its Gaussian density energy obeys
\begin{align}
 E_{H,N}
 &=\frac{dk_1}{N}\|V^1\|_F^2
   +\sum_{l=2}^H\frac{k_l}{k_{l-1}}\|V^l\|_F^2
   +\frac N{k_H}\|b\|^2\\
 &\le d(d+1)B^2p_1+\alpha(H-1)C^2B^2
                        +\frac{2C^2B^2}p.
 \label{eq:geometric-mean-energy}
\end{align}
The geometric increase pays a constant energy per internal layer, rather than a factor growing with learner width.

Let $\nu$ be a Gaussian law centered at this mean network. Multiply all first-layer prior variances by $p_1$, and leave every internal and readout variance unchanged. Exact Gaussian relative entropy gives
\begin{equation}
 \frac2\omega\KL(\nu\Vert\mu)
 =E_{H,N}+\tau^2(d+1)(p_1-1-\log p_1)
 \le E_{H,N}+(d+1)[\log(2/p)+(H-1)\log\alpha]
 \le c_1H.
 \label{eq:geometric-kl}
\end{equation}
Here and below constants may depend on $d,C,B,p$. The only variance suppression occurs in the first layer. Its logarithmic cost is proportional to depth.

\subsection{The native Gaussian bulk remains small}
Fix an input. Let $q_l=N^{-1}\E_\nu\|h^l\|^2$, and let $e_l$ be the largest expected squared error of an active clone relative to its comparison neuron's hidden value. The first layer satisfies
\[
 q_1\le p_1+2p_1=3p_1,\qquad e_1\le2p_1.
\]
The active contribution to $q_l$ is at most $p_l$, because the activation is bounded by one. Each inactive current row has zero mean. Conditional on the preceding layer its preactivation variance is $\tau^2\|h^{l-1}\|^2/N$. Since $|\phi(t)|\le|t|$, it follows that
\[
 q_l\le p_l+\tau^2q_{l-1}\le p_l+q_{l-1}\le3p_l.
 \label{eq:geometric-q}
\]
The final inequality follows inductively from $p_l=\alpha p_{l-1}$ and $\alpha\ge2$.

For an active row, the deterministic part of its preactivation error is a weighted sum of preceding clone errors whose absolute coefficient sum is at most $A$. Its $L^2(\nu)$ norm is therefore at most $A\sqrt{e_{l-1}}$. The current centered Gaussian row noise is independent of the preceding layer. Its cross term with that deterministic error has expectation zero and its expected variance is $\tau^2q_{l-1}$. Applying the activation's Lipschitz bound gives
\[
 e_l\le A^2e_{l-1}+\tau^2q_{l-1}
     \le A^2e_{l-1}+3p_{l-1}\le2p_l.
 \label{eq:geometric-e}
\]
No independence between clones is assumed. The last inequality uses the chosen $\alpha$.

The deterministic readout coefficients have total absolute sum at most $A$, so its expected squared error is at most $A^2e_H$. The independent centered readout noise adds variance $\tau^2q_H/N$. Thus, uniformly over the input cube,
\begin{equation}
 \E_\nu|z_w(x)-g(x)|^2
 \le A^2e_H+\tau^2q_H/N
 \le(2A^2+3)p.
 \label{eq:geometric-output-error}
\end{equation}
Choose $p$ sufficiently small that the square root of the last quantity is at most $\epsilon/[128(CB+1)]$. Since $|g|\le CB$ and labels are bounded by one, this makes the comparison's excess empirical loss at most $\epsilon/32$ for every sample. Combining \eqref{eq:geometric-kl} with the posterior variational identity and choosing $\ell_+$ sufficiently small yields
\[
 \int\widehat R_S(z_w)d\Pi\le\widehat R_S(g)+\epsilon/16
\]
at every width above the stated exponential threshold. Below it, the ordinary clone neighborhood of Lemma~\ref{lem:teacher} gives normalized entropy $O(H^2)$; choosing $\ell_-$ small yields the same conclusion with $\lambda=\ell_-/H^2$.

\subsection{A common finite-cover argument for all widths and scales}
Use the total operator energy
\[
 Q_{\mathrm{op}}=\|W^1\|_F^2/N+
       \sum_{l=2}^H\|W^l\|_\op^2+N\|v\|^2.
\]
Its Gaussian moment estimate \eqref{eq:relu-total-op-mgf} depends only on the weight law and applies here. The deterministic evidence bound $Z\ge e^{-\omega/\lambda}$ and prior bound $\E_\mu R(z)^2\le24$ imply that the posterior raw-risk contribution outside
\[
 Q_{\mathrm{op}}\le T_0H\quad(N\ge N_*),\qquad
 Q_{\mathrm{op}}\le T_0H^2\quad(N<N_*)
\]
is at most $\epsilon/64$ for sufficiently large fixed $T_0$. Indeed Cauchy--Schwarz bounds that contribution by
\[
 \sqrt{24}\exp[\omega/\lambda+C_0N(H+1)/2-c_0\omega Q_*/2],
\]
where $Q_*$ is the appropriate cutoff. Use $N\le\omega$ and the two specified likelihood scales.

At large widths the raw output amplitude is at most $\sqrt{T_0H}$. AM--GM bounds its input Lipschitz constant by
$(T_0H/(H+1))^{(H+1)/2}\le e^{c_2H}$. A spatial cover at fixed dimension has log cardinality $e^{O(H)}$. Its square-loss envelope is polynomial in $H$, so finite-class Hoeffding still costs only $e^{O(H)}$ samples.

At smaller widths, $N<e^{aH}$ and the good-set raw amplitude is at most $\sqrt{T_0}H$. There are $e^{O(H)}$ parameters. Their coordinate magnitudes are at most $\sqrt{NT_0H^2}$, and a telescoping bound gives a parameter-to-output Lipschitz constant whose logarithm is $O(H^2)$. A union of finite parameter grids over these widths therefore has log cardinality $e^{O(H)}$. Its loss envelope is again polynomial in $H$. Both dictionaries are independent of $\tau$ and of the sample.

Choose their uniform output accuracy small enough that loss perturbations on the stated amplitude intervals are at most $\epsilon/64$; accuracy $c_\epsilon/H$ suffices for both. This only adds logarithmic factors to the cover sizes. Finite-class concentration gives one event on which every good network has population loss at most its empirical loss plus $\epsilon/8$, for all widths and scales. Integrate, use the empirical objective bound and the raw-risk tail, and apply scalar concentration to a fixed population minimizer in the compact comparison class. Increasing $c$ absorbs all constants. This proves Theorem~\ref{thm:geometric-depth}, and Jensen gives its ordinary posterior-mean conclusion.


\section{A depth upper bound that includes mean field}
\label{sec:all-scales}
The preceding exponential-depth proof uses $\tau\to0$ with width. A different finite-cover argument gives a weaker depth bound at every $0<\tau\le1$, including canonical mean field $\tau=1$. Here set $\omega=N/\tau^2$ without imposing $\tau^2=N^{1-\gamma}$.

\begin{theorem}[All scales, with explicit dimension dependence]
\label{thm:all-scales}
Fix $C,B,\epsilon,\delta$ as before and let $T=d+H+2$. There are $c>0$ and $0<\ell<1$, independent of $d,H,N,\tau$, such that the posterior with $\lambda=\ell/T^2$ satisfies \eqref{eq:main}, on one sample event common to every $N\ge C$ and every $0<\tau\le1$, whenever
\begin{equation}
 n\ge T^{c(d+1)H}.
 \label{eq:spatial-depth-budget}
\end{equation}
If $\phi$ is twice continuously differentiable with a fixed bound on $|\phi''|$, a second sufficient budget is
\begin{equation}
 n\ge T^{cH^2},
 \label{eq:program-depth-budget}
\end{equation}
with $c$ also depending on that derivative bound. Thus the smooth bounded case permits $n\ge T^{cH\min\{H,d+1\}}$ after enlarging $c$.
\end{theorem}

At fixed dimension the spatial bound is $\exp(O(H\log H))$. It has an exponential dependence on dimension. The program bound is polynomial in dimension at every fixed depth; for $d\le H$ it is $\exp(O(H^2\log H))$. These distinct dimension dependences should not be conflated.

\subsection{A common comparison and tail estimate}
Use Lemma~\ref{lem:teacher} with arbitrary $\tau\le1$. Its clone energy is $O_{C,B}(T^2)$, cylinder dimension is $O_C(T)$, and $\log P_H=O_{C,B}(T)$. With $\lambda=\ell/T^2$, its comparison cost obeys
\[
 A_{H,N}(\lambda)\le c_{C,B}T^2[1+\log(1/\ell)].
\]
Choose $\ell$ small enough to make the comparison penalty a prescribed fraction of $\epsilon$.

For an absolute sufficiently large lower cutoff, the norm event
\begin{equation}
 \|W^1\|_F^2/N\le K,\qquad
 \max_{2\le l\le H}\|W^l\|_\op^2\le K,\qquad
 N\|v\|^2\le K
 \label{eq:all-scale-norm}
\end{equation}
has prior complement at most $c(H+1)e^{-c_0\omega K}$, with absolute $c,c_0>0$. This follows from chi-square and Gaussian operator-norm tails. The factor $1/d$ in the first-layer variance keeps its mean energy at most two, independently of dimension.
Choose $K=k_0T^2$, with fixed $k_0$ sufficiently large. A parameter-good set with the same prior complement bound has posterior raw-risk contribution on its complement at most
\begin{equation}
 \sqrt{24c(H+1)}\exp[\omega/\lambda-c_0\omega K/2]
 \le\epsilon/32.
 \label{eq:all-scale-bad-risk}
\end{equation}
This uses the evidence lower bound and prior risk second moment already proved. The same inequality is valid for a good set depending on a fixed sample, provided its prior bound holds for each such sample. The high-empirical-risk tail on $\widehat R>8$ is at most $2e^{-3\omega/\lambda}$.

It remains to prove $R(z)\le\widehat R(z)+\epsilon/4$ on the relevant good set with empirical risk at most eight. Integrating, using the teacher objective, and applying one fixed minimizer's scalar concentration event then proves the theorem.

\subsection{The spatial bound}
On \eqref{eq:all-scale-norm}, the raw output has amplitude at most $M=\sqrt K$ and Euclidean input Lipschitz constant at most
$L_K=K^{(H+1)/2}$; the first-layer $\sqrt N$ cancels the readout's $N^{-1/2}$.
A coordinate grid of spacing $a/(4L_K\sqrt d)$ and grid-value quantization at spacing $a/4$ gives a finite uniform-norm $a$-cover of this entire function class. Nearest-grid representatives are measurable piecewise constant functions. The dictionary $\mathcal D$ satisfies
\[
 \log|\mathcal D|\le(c\sqrt d L_K/a)^d\log(cM/a).
\]
Choose fixed $a=c_\epsilon\epsilon$ sufficiently small. Finite-class Hoeffding gives population-minus-empirical loss at most $\epsilon/32$ for every dictionary member if
\[
 n\ge c\epsilon^{-2}(M+1)^4
                 [\log|\mathcal D|+\log(c/\delta)].
\]
On networks with empirical risk at most eight, the approximant has bounded empirical and population risks. Minkowski's inequality for both square-loss norms transfers the bound with $O(a)$ loss cost. Since $K=k_0T^2$, this sample condition is implied by \eqref{eq:spatial-depth-budget}. The dictionary is independent of width, scale, and the data law, so its statistical event is common to all learners.

\subsection{A Gaussian finite dictionary with tracked constants}
We provide the additional compression argument for \eqref{eq:program-depth-budget}. Constants in this subsection can depend on the fixed derivative bound, but their dependence on depth is tracked.
Fix any probability measure $\xi$ on the input cube. For an actual hidden layer define the normalized covariance operator
\[
 Q_l(x,y)=N^{-1}\sum_i h_i^l(x)h_i^l(y)
\]
on $L^2(\xi)$. Its trace is at most one. At cutoff $\eta_l$, retain factors $F_l=(\sqrt{\lambda_j}e_j)_j$ of rank $r_l\le1/\eta_l$ and let
$v_l(x)=Q_l(x,x)-\|F_l(x)\|^2\in[0,1]$.
There are exact row coefficients with
\[
 F_l=\E_i b_i h_i^l,\qquad \E_i b_i b_i^\top=I,
\]
where $\E_i$ is the uniform finite row average. This follows from the orthonormal right-singular vectors of the normalized hidden-feature operator.

Conditional on the preceding layer, a current Gaussian row field is
$a_i\cdot F_{l-1}+Z_i^0$, where $a_i$ is centered Gaussian with covariance $\tau^2I$ and the independent residual field has covariance $\tau^2E$, with $\|E\|_\op\le\eta_{l-1}$ and diagonal at most one. Write $r_i$ for the activated feature minus its conditional mean given $a_i$. Uniformly over every coefficient vector,
\begin{equation}
 \left\|c\mapsto N^{-1/2}\sum_i c_i r_i\right\|_\op
 \le c[\omega^{-1/2}+\sqrt{K\eta_{l-1}}]
 \label{eq:depth-residual}
\end{equation}
except on conditional probability $e^{-\omega K}$. To prove this, fix $\|c\|_2=1$. Gaussian Poincar\'e gives mean Hilbert norm at most $\tau/\sqrt N$, and Gaussian concentration has Lipschitz constant $\tau\sqrt{\eta_{l-1}/N}$. A $1/2$-net of the coefficient sphere has at most $5^N$ points; its cost is absorbed into the displayed deviation because $\tau\le1$. The same argument for $\phi^2$, which is two-Lipschitz, bounds the centered diagonal-profile mean without the coefficient net. The operator bound holds also for coefficients selected from the current realization.

Let $\psi(m,v)=\E\phi(m+\sqrt vG)$ and $\chi(m,v)=\E\phi(m+\sqrt vG)^2$. The Gaussian heat equation makes both globally Lipschitz in $m$ and in $v\in[0,1]$, with fixed constants. For $\psi_i=\psi(a_i\cdot F_{l-1},\tau^2v_{l-1})$ and $\chi_i=\chi(a_i\cdot F_{l-1},\tau^2v_{l-1})$, the exact conditional-mean factor and residual variance are
\[
 F_l^0=\E_i b_i\psi_i,\qquad
 v_l^0=\E_i\chi_i-\|F_l^0\|^2
      =\E_i\chi_i-\E_{i,j}(b_i\cdot b_j)\psi_i\psi_j.
\]
Bessel's inequality gives $\|F_l^0(x)\|^2\le\E_i\psi_i(x)^2\le\E_i\chi_i(x)\le1$ and $v_l^0\in[0,1]$. On the norm event, the exact latent Gram satisfies $\E_i a_i a_i^\top\preceq KI$. Thus preceding factor and variance errors propagate through these means by at most
$c_\phi[\sqrt K e_{l-1}^F+e_{l-1}^V]$, without a rank multiplier. Deleting the residual fields costs at most
$c\sqrt{r_l+1}[\omega^{-1/2}+\sqrt{K\eta_{l-1}}]$ for both outputs; use \eqref{eq:depth-residual} for the factors and the diagonal bound for the variance.

Approximate these finite means by sampling with replacement after evaluating on the already chosen preceding approximants. Factor integrands have Hilbert second moment at most $r_l$, while the scalar pair coefficient satisfies
$\E_{i,j}(b_i\cdot b_j)^2=r_l$.
Hence sample counts $cH(r_l+1)/(\delta_l^F)^2$ and $cH(r_l+1)/(\delta_l^V)^2$ give the respective errors with positive probability. Clip only variance approximants to $[0,1]$; leave factor approximants linear. Their exact means remain pointwise bounded even when the preceding sampled factor vector is not. The first layer is sampled directly from its actual features and needs no residual deletion or minimum width. The final output is approximated directly as a scalar average with coefficient second moment at most $K$, using $cHK/a^2$ samples.

Here is an explicit error schedule. Put $G_*=c_\phi\sqrt K$ with a sufficiently large fixed $c_\phi$, and for $t=H-l$ choose
\[
 \delta_l^F=\frac{a}{cH G_*^{t+1}},\qquad
 \delta_l^V=\sqrt K\,\delta_l^F,\qquad
 \eta_{H-1}=\frac{a^2}{cH^2K^2},\qquad
 \eta_{l-1}=\frac{(\delta_l^F)^2\eta_l}{c_\phi K}.
\]
These choices make every layer's propagated contribution at most $a/(cH)$. The factor error is smaller than its tolerance when the effective width exceeds
\[
 \omega_* =cH\max\{K/a^2, \max_{2\le l\le H-1}(r_l+1)/(\delta_l^F)^2\}.
\]
The inner maximum is omitted when $H=2$. For fixed $a$, $K=k_0T^2$, and $H\le T$, this schedule gives
\[
 \log(1/\eta_l)\le c_\phi H^2\log T,\qquad
 \log(1/\delta_l^F)\le c_\phi H\log T,\qquad
 \omega_*\le T^{c_\phi H^2}.
\]
All ranks and sample counts are therefore at most $T^{c_\phi H^2}$.

Markov bounds control the sampled coefficients; the logarithms of their caps are $O_\phi(H^2\log T)$. Linear sampled factors can be contracted into subsequent incoming coefficients. Variances use the separate scalar pair representation, so they do not require the norm of a sampled factor vector. The programs form a shared layered directed graph: the parameter count sums adjacent-layer sample-count products, rather than unfolding a separate tree for every output. It is at most $T^{c_\phi H^2}$. Parameter magnitudes obey the same bound. With $O(H)$ computational stages of bounded Gaussian means and scalar clipping, the parameter-to-uniform-function Lipschitz constant has logarithm at most $c_\phi H^3\log T$. This larger quantity enters only the logarithm of a parameter grid. Consequently a fixed finite dictionary of accuracy $a$ has
\begin{equation}
 \log|\mathcal D|\le T^{c_\phi H^2}.
 \label{eq:depth-program-entropy}
\end{equation}
Clip final program outputs to $[-\sqrt K,\sqrt K]$.

If $\omega<\omega_*$, then $N\le\omega<\omega_*$. Include actual norm-bounded networks at every such integer width in the same dictionary. Their parameter counts are $O_H(dN+N^2)\le T^{c_\phi H^2}$, with polynomially bounded weights and the same bound on the logarithm of their parameter-to-function Lipschitz constants. Uniform parameter grids and a finite union preserve \eqref{eq:depth-program-entropy}.

This constructs, for each fixed $\xi$, a prior-good set of complement at most $c(H+1)e^{-c_0\omega K}$ whose outputs have $L^2(\xi)$ approximants of error $a$ in one fixed dictionary. The dictionary is independent of $\xi,N,\tau$ and the sample. Good sets are Borel, using measurable spectral bases of finite Gram matrices; all probability estimates follow from conditional Gaussian bounds and a union over layers.

Use $\xi=(D_X+\widehat D_{X,S})/2$. Then each approximant is close in both population and empirical square-loss norms. The auxiliary good set may depend on the input sample, but its prior complement bound holds for every fixed sample, so \eqref{eq:all-scale-bad-risk} applies pointwise. The statistical event is applied only to the fixed dictionary. Finite-class Hoeffding costs a polynomial in $K$ times \eqref{eq:depth-program-entropy}, which is absorbed into \eqref{eq:program-depth-budget}. Root-risk transfer on empirical risk at most eight gives the required deviation. This completes the all-scale theorem.


\section{An all-scale exponential-depth bound for saturated ReLU prediction}
\label{sec:relu-all-scale-depth}
The preceding raw-output results assume bounded hidden activations. ReLU admits a different construction because its positive homogeneity lets a small active subnetwork carry signals of size $\sqrt N$. This yields an exponential-depth theorem even at canonical mean field, provided the observed output is saturated. The saturation is part of both the likelihood and the prediction rule.

Let $\phi(t)=\max\{t,0\}$ in every hidden layer, and put $s(t)=\max\{-1,\min\{1,t\}\}$. The observed network output is $f_w=s(z_w)$. Use the same independent Gaussian variances as before, with any $0<\tau\le1$ and $\omega=N/\tau^2$, and define
\[
 d\Pi_{N,H,S}(w)=Z^{-1}\exp[-\omega\widehat R_S(f_w)/(2\lambda)]d\mu_{N,H}(w).
\]
The prediction is $\bar f(x)=\int s(z_w(x))d\Pi(w)$. In particular it is not $s(\int z_wd\Pi)$, and this section does not use a likelihood on the raw outputs.
Let $\mathcal G^{\mathrm{sat}}_{C,B,H}$ be the class of saturated, same-depth, width-$C$ ReLU comparison networks with every coefficient bounded by $B$.

\begin{theorem}[All scales, saturated ReLU]
\label{thm:relu-all-scale-depth}
Fix $d,C,B,\epsilon,\delta$ as in Theorem~\ref{thm:main}. There are constants $a,c>0$ and $0<\ell_+,\ell_-<1$, depending only on those fixed resources, such that the following holds for every $H\ge2$. Set
\[
 N_*(H)=\lceil e^{aH}\rceil,\qquad
 \lambda_{H,N}=
 \begin{cases}\ell_-/H^2,&N<N_*(H),\\
 \ell_+/H,&N\ge N_*(H).
 \end{cases}
\]
If $n\ge e^{cH}$, one event of sample probability at least $1-\delta$ satisfies, simultaneously for every $N\ge C$ and every $0<\tau\le1$,
\[
 \int R(f_w)d\Pi_{N,H,S}
 \le\inf_{g\in\mathcal G^{\mathrm{sat}}_{C,B,H}}R(g)+\epsilon.
\]
The same bound holds for $R(\bar f)$. Thus every fixed power-law prior exponent $\gamma\ge1$, including $\gamma=1$, is covered.
\end{theorem}

The same proof applies if $s$ is replaced by any fixed one-Lipschitz map $s:\mathbb R\to[-1,1]$, including a tanh readout, used in both the learner and comparison class. It also applies to positively homogeneous one-Lipschitz hidden activations such as leaky ReLU.

The constants are not uniform in dimension or comparison resources. This is a fixed-resource depth theorem. It does not assert that the canonical raw-output tanh learner has the same bound.

\subsection{A translated Gaussian comparison with no variance shrinkage}
Write a comparison network's hidden matrices and readout as $V^1,\ldots,V^H,b$. Choose $C$ active neurons in every learner layer. The deterministic mean network has first active rows $\sqrt N V^1$, internal active blocks $V^l$, readout $b/\sqrt N$ on active neurons, and zero mean elsewhere. ReLU homogeneity makes its active hidden vectors exactly $\sqrt N$ times the comparison vectors, and its raw output is the comparison output.
Its density energy is
\begin{equation}
 E_H=d\|V^1\|_F^2+\sum_{l=2}^H\|V^l\|_F^2+\|b\|^2
 \le dC(d+1)B^2+(H-1)C^2B^2+CB^2
 \le c_1H.
 \label{eq:relu-mean-energy}
\end{equation}
For $0<r\le1$, center a Gaussian comparison law $\nu_r$ at this mean. Multiply prior variances by $r^2$ only on the active first-layer rows, the rows entering active internal neurons, and the full readout. All remaining coordinates keep their prior variances. The number of changed variances is
\[
 D_{H,N}=C(d+1)+(H-1)CN+N.
\]
The exact Gaussian relative entropy gives
\begin{equation}
 \frac{2}{\omega}\KL(\nu_r\Vert\mu)
 =E_H+\frac{D_{H,N}}{\omega}(r^2-1-2\log r)
 \le E_H+2\frac{D_{H,N}}N\log(1/r).
 \label{eq:relu-selective-entropy-depth}
\end{equation}
In particular $r=1$ translates the full prior without changing any covariance and pays precisely $E_H$.

We next keep the width factor in its approximation error. For a fixed input, let $q_l=N^{-1}\E_{\nu_r}\|h^l\|^2$, and let $e_l$ be $N^{-1}$ times the expected squared difference between the active hidden vector and its deterministic comparison value. Put $A=CB$ and $K_0=A^2+1$. There is $Q_0$, depending only on $d,C,B$, such that
\[
 q_1\le Q_0,\qquad e_1\le 2Cr^2/N,
\]
and, for $l\ge2$,
\begin{equation}
 q_l\le K_0 q_{l-1},\qquad
 e_l\le A^2e_{l-1}+\frac{Cr^2}{N}q_{l-1}.
 \label{eq:relu-depth-moments}
\end{equation}
Indeed ReLU has norm at most that of its preactivation and is one-Lipschitz. The current centered Gaussian row noise is independent of the preceding hidden vector, so all mean--noise cross terms vanish before the norm bounds are applied. The active noise contributes $Cr^2\tau^2q_{l-1}/N$ to normalized squared error; the full noise contributes at most $\tau^2q_{l-1}$ to $q_l$. We have used $\tau\le1$.

Solving \eqref{eq:relu-depth-moments} and then integrating the independent readout noise gives
\begin{equation}
 \sup_x\E_{\nu_r}|z_w(x)-z_g(x)|^2
 \le \|b\|^2 e_H+\frac{r^2}{N}q_H
 \le\frac{r^2J_H}{N},\qquad
 J_H\le c_2(H+1)K_0^H\le e^{c_3H}.
 \label{eq:relu-unshrunk-depth-error}
\end{equation}
The supremum is over the input cube. Since $s$ is one-Lipschitz, the same bound holds for saturated output differences. For labels in $[-1,1]$, squared loss is four-Lipschitz on this output interval. Hence, for every fixed sample,
\[
 \int\widehat R_S(f_w)d\nu_r
 \le\widehat R_S(g)+4r\sqrt{J_H/N}.
\]
Choose $N_*(H)$ so large that $4\sqrt{J_H/N_*}\le\epsilon/64$; an exponential $e^{aH}$ suffices. At larger widths use $r=1$. At smaller widths use
$r_H=\min\{1,\epsilon/(256\sqrt{J_H})\}$.
Then both regimes have comparison loss error at most $\epsilon/64$. Equations \eqref{eq:relu-mean-energy}--\eqref{eq:relu-selective-entropy-depth} give normalized entropy costs $O(H)$ and $O(H^2)$ respectively, uniformly over $N$ and $\tau$. Choose $\ell_+,\ell_-$ sufficiently small. The posterior variational identity then yields, in both regimes,
\begin{equation}
 \int\widehat R_S(f_w)d\Pi
 \le\widehat R_S(g)+\epsilon/16.
 \label{eq:relu-depth-objective}
\end{equation}
The full variational objective also obeys this bound; we have dropped its nonnegative entropy term.

\subsection{A total operator-energy bound}
Define
\[
 Q_{\mathrm{op}}(w)=\|W^1\|_F^2/N+
       \sum_{l=2}^H\|W^l\|_\op^2+N\|v\|^2.
\]
For absolute constants $c_0,C_0>0$,
\begin{equation}
 \E_\mu e^{c_0\omega Q_{\mathrm{op}}}
 \le e^{C_0N(H+1)}.
 \label{eq:relu-total-op-mgf}
\end{equation}
The first and last factors follow directly from chi-square moment generating functions. For an internal standard Gaussian matrix $G$, the bound
$\Pr(\|G\|_\op>2\sqrt N+t)\le e^{-t^2/2}$ gives
$\E e^{c_0\|G\|_\op^2}\le e^{C_0N}$ for sufficiently small $c_0$.
The scaling $W=\tau G/\sqrt N$ cancels exactly against $\omega=N/\tau^2$. Independence multiplies the moment bounds.

Every saturated loss is at most four, so $Z\ge e^{-2\omega/\lambda}$. Since $N\le\omega$, \eqref{eq:relu-total-op-mgf} implies that, for a sufficiently large fixed $T_0$, the posterior mass outside
\[
 Q_{\mathrm{op}}\le T_0H\quad(N\ge N_*),\qquad
 Q_{\mathrm{op}}\le T_0H^2\quad(N<N_*)
\]
is at most $\epsilon/64$ in both regimes. These bounds hold for every sample. No posterior concentration limit is being taken.

\subsection{Two finite covers and one common sample event}
For the large-width good set, the input Lipschitz constant of $f_w$ is at most
\[
 \sqrt{N\|v\|^2}\sqrt{\|W^1\|_F^2/N}
       \prod_{l=2}^H\|W^l\|_\op
 \le\left(\frac{Q_{\mathrm{op}}}{H+1}\right)^{(H+1)/2}
 \le e^{c_4H}.
\]
The factors $N^{1/2}$ and $N^{-1/2}$ cancel, and the middle inequality is AM--GM. All outputs are in $[-1,1]$. A spatial grid and quantized values therefore give one fixed uniform-norm cover of accuracy $\epsilon/256$ whose log cardinality is at most $e^{c_5H}$ at fixed $d$.

For the small-width good set, include every integer $C\le N<N_*$. There are at most
$O(HN_*^2+dN_*)=e^{O(H)}$ parameters per network. The energy bound gives coordinate magnitudes at most $\sqrt{NT_0H^2}$. A layerwise telescoping estimate makes the logarithm of the parameter-to-output Lipschitz constant $O(H\log(N+H+2))=O(H^2)$. Uniform parameter grids, followed by final saturation, thus give a second fixed cover of accuracy $\epsilon/256$ with log cardinality at most $e^{c_6H}$. Taking the union over these widths changes only the constant. The covers do not depend on the scale $\tau$.

Finite-class Hoeffding on the union of both dictionaries gives a common event on which every good function satisfies
$R(f)\le\widehat R(f)+\epsilon/8$, once $n\ge e^{cH}$. Loss perturbations from the uniform covers cost at most four times their output accuracy. Integrating this inequality, charging at most four times the posterior bad-set mass, and using \eqref{eq:relu-depth-objective} gives
\[
 \int R(f_w)d\Pi\le\widehat R(g)+\epsilon/4.
\]
The comparison parameter cube is compact and its loss is continuous, so a population-risk minimizer exists. Apply scalar concentration only to this fixed minimizer; its outputs and labels are bounded. Enlarge $c$ to absorb that event and all fixed constants. This proves Theorem~\ref{thm:relu-all-scale-depth} simultaneously for every stated width and scale. Jensen's inequality gives the posterior-mean conclusion.


\subsection{A depth-dependent strengthening of mean field}
\label{sec:depth-normalized}
A small depth-dependent prior prefactor also gives an exponential-depth bound without restricting the support of the Gaussian prior.

\begin{corollary}[Bounded hidden activations with $H\tau^2\le1$]
\label{cor:depth-normalized}
Fix $d,C,B,\epsilon,\delta$, and let the hidden activation be any one-Lipschitz map into $[-1,1]$. There are $c>0$ and $\ell>0$, depending only on those fixed resources, such that $n\ge e^{cH}$ and $\lambda=\ell/H$ imply \eqref{eq:main} simultaneously for every $N\ge C$ and every $0<\tau^2\le1/H$. This uses raw Gaussian likelihood and raw-output prediction, as in Theorem~\ref{thm:main}.
\end{corollary}

\begin{proof}
The exact comparison cost \eqref{eq:exact-comparison-cost}, with $\lambda=\ell/H$, is at most
$c_1H+c_2\tau^2H^2[1+\log(1/\ell)]\le c_3H[1+\log(1/\ell)]$.
Choose $\ell$ sufficiently small to make the comparison penalty a prescribed fraction of $\epsilon$.
The moment estimate \eqref{eq:relu-total-op-mgf} concerns only Gaussian matrices and therefore applies unchanged to bounded activations. Choose a sufficiently large fixed $T_0$. Since $N/\omega=\tau^2\le1/H$, the prior mass of $Q_{\mathrm{op}}>T_0H$ is at most
$\exp[C_0N(H+1)-c_0\omega T_0H]$.
The evidence and raw risk second-moment bounds from the bounded-activation proof make the posterior raw-risk contribution of this bad set arbitrarily small, uniformly over all stated widths and scales.

On the good set, the raw output has magnitude at most $\sqrt{T_0H}$, while its input Lipschitz constant is at most
$(T_0H/(H+1))^{(H+1)/2}\le e^{c_4H}$.
The amplitude estimate uses bounded hidden features, and the Lipschitz estimate uses the product of operator norms and AM--GM. The spatial-cover proof of Section~\ref{sec:all-scales} now has log dictionary size $e^{O(H)}$ at fixed $d$. Its polynomial loss-envelope factor is absorbed into this exponential. Apply its statistical event, posterior tail integration and one fixed comparison minimizer's concentration event. These events depend only on the fixed dictionary, so they cover all $N$ and $\tau$ simultaneously.
\end{proof}

For a power law $\tau^2=N^{1-\gamma}$ with fixed $\gamma>1$, the corollary applies once $N\ge H^{1/(\gamma-1)}$. Below that polynomial width threshold, the finite-width PAC--Bayes estimate \eqref{eq:finite-width} with $\lambda=\ell_-/H^2$ suffices using at most exponentially many samples. Thus, at fixed input dimension, the exponential width threshold in Theorem~\ref{thm:main} may be replaced by this polynomial threshold. Its stronger growing-dimension conclusion uses the spectral proof instead.

The condition $\tau^2\le1/H$ changes the prior variance as depth changes. At $\tau^2=1/H$, the internal variances are $1/(HN)$, while canonical mean field has $1/N$. The corollary therefore does not settle the canonical raw-output bounded-activation case.


The proof has undergone independent AI-agent mathematical review and is not Lean formalized.
\begin{thebibliography}{9}
\bibitem{talagrand} Michel Talagrand. \emph{Upper and Lower Bounds for Stochastic Processes}. Springer, 2014.
\bibitem{sst} Nathan Srebro, Karthik Sridharan, and Ambuj Tewari. \emph{Smoothness, Low Noise and Fast Rates}. Advances in Neural Information Processing Systems 23, 2010. \url{https://www.cs.cornell.edu/~sridharan/smoothness.pdf}.
\end{thebibliography}
\end{document}
